% META: {"id": "TSPE-LIMFCT-EV-A", "chapitre": "TSPE-LIMITES-FONCTIONS", "type_objet": "evaluation", "version": "A", "duree_min": 55, "bareme_total": 20, "capacites_codes": ["C1", "C2", "C3"], "competences": ["chercher", "calculer", "raisonner", "communiquer"], "mode_creation": "ex_nihilo", "sources_inspiration": [], "status": "generated"}
% BEGIN-VERIFY
% from sympy import *
% x = symbols('x')
%
% # ---- Exercice 1 (6 pts) : Limites (C1) ----
% assert limit((3*x**2 + x - 1)/(x**2 - 2), x, oo) == 3
% assert limit(5*Rational(1,3)**x + 2, x, oo) == 2
% assert limit(x**2 - 4*x + 1, x, oo) == oo
%
% # ---- Exercice 2 (7 pts) : Asymptotes (C2) ----
% f2 = (2*x + 5)/(x - 1)
% assert limit(f2, x, oo) == 2
% assert limit(f2, x, 1, '+') == oo
% assert limit(f2, x, 1, '-') == -oo
% assert simplify(f2 - (2 + 7/(x-1))) == 0
%
% # ---- Exercice 3 (7 pts) : Croissances comparees (C3) ----
% f3 = (x**2 - 3)*exp(-x)
% assert limit(f3, x, oo) == 0
% assert limit(f3, x, -oo) == oo
% f3p = diff(f3, x)
% assert simplify(f3p - (-x**2 + 2*x + 3)*exp(-x)) == 0
% assert simplify(f3p - (-(x - 3)*(x + 1))*exp(-x)) == 0
% assert f3.subs(x, 3) == 6*exp(-3)
% END-VERIFY

%---------------------------------------------------------------
% Duree : 55 minutes — Bareme : 20 points
% Toute reponse doit etre justifiee ; les resultats seuls ne sont pas acceptes.
%---------------------------------------------------------------

\begin{evaluation}{TSPE-LIMFCT-EV-A}{55}

%===============================================================
% EXERCICE 1 — Limites de fonctions (6 points)
%===============================================================

\begin{exercice}{TSPE-LIMFCT-EV-A-EX1}{2}{12}

\textbf{Exercice 1} (6 points) — \textit{Capacite C1}

\medskip

Determiner les limites suivantes en justifiant.

\begin{enumerate}
  \item (2 pts) $\lim_{x \to +\infty} \dfrac{3x^2 + x - 1}{x^2 - 2}$
  \item (2 pts) $\lim_{x \to +\infty} \left(5 \times \left(\dfrac{1}{3}\right)^x + 2\right)$
  \item (2 pts) $\lim_{x \to +\infty} (x^2 - 4x + 1)$ (factoriser par le terme preponderant)
\end{enumerate}

\end{exercice}

%===============================================================
% EXERCICE 2 — Asymptotes (7 points)
%===============================================================

\begin{exercice}{TSPE-LIMFCT-EV-A-EX2}{2}{15}

\textbf{Exercice 2} (7 points) — \textit{Capacite C2}

\medskip

Soit $f(x) = \dfrac{2x + 5}{x - 1}$ definie sur $\mathbb{R} \setminus \{1\}$.

\begin{enumerate}
  \item (2 pts) Determiner $\lim_{x \to +\infty} f(x)$ et $\lim_{x \to -\infty} f(x)$. En deduire l'asymptote horizontale.
  \item (2 pts) Determiner $\lim_{x \to 1^+} f(x)$ et $\lim_{x \to 1^-} f(x)$. En deduire l'asymptote verticale.
  \item (1,5 pt) Montrer que $f(x) = 2 + \dfrac{7}{x-1}$.
  \item (1,5 pt) En deduire la position de la courbe par rapport a l'asymptote horizontale.
\end{enumerate}

\end{exercice}

%===============================================================
% EXERCICE 3 — Etude de fonction avec exponentielle (7 points)
%===============================================================

\begin{exercice}{TSPE-LIMFCT-EV-A-EX3}{3}{20}

\textbf{Exercice 3} (7 points) — \textit{Capacite C3}

\medskip

Soit $f(x) = (x^2 - 3)\,\mathrm{e}^{-x}$ definie sur $\mathbb{R}$.

\begin{enumerate}
  \item (2 pts) Determiner $\lim_{x \to +\infty} f(x)$ et $\lim_{x \to -\infty} f(x)$. Preciser l'asymptote eventuelle.
  \item (2,5 pts) Calculer $f'(x)$ et montrer que $f'(x) = -(x-3)(x+1)\,\mathrm{e}^{-x}$.
  \item (1,5 pt) Dresser le tableau de variations de $f$.
  \item (1 pt) Calculer $f(3)$.
\end{enumerate}

\end{exercice}

\end{evaluation}
